%=============================================================================
\section{Adoption-risk scenarios for outcomes not measured in Spain}
\label{sec:scenarios}
%=============================================================================

This section concerns the five outcomes absent from the CIS; only one of them
is projected. The projections are prior-informed conditional estimates: their
levels and their profile
effects come from external evidence and from the assumptions of
Eq.~\eqref{eq:transfer}. They are not observed Spanish values and none has been
validated in Spain.

%-----------------------------------------------------------------------------
\subsection{Evidence base for the target outcomes}
\label{sec:evidence_base}
%-----------------------------------------------------------------------------

Table~\ref{tab:novel} summarises the pooled evidence for each target outcome.
The base is thin. General acceptance rests on two studies; food-packaging
acceptance on one item from one Finnish survey, in which 38\% of respondents
judged recycled plastic suitable for food packaging \citep{ruokamo2022};
perceived safety on a single sample of 116 respondents; perceived quality on
no usable entry;
and willingness to pay a premium on two studies, one of which concerns
bioplastics.

\begin{table}[htbp]
\centering
\caption{Pooled external evidence for the five target outcomes. Low, middle,
and high are the two lowest, the central, and the two highest of five response
categories. Food-packaging acceptance is reported as its source statistic.}
\label{tab:novel}
\small
\begin{threeparttable}
\begin{tabular}{lrrrccc}
\toprule
\textbf{Outcome} & \textbf{Studies} & \textbf{Entries} &
  \textbf{Indep.\ $n$} & \textbf{Low} & \textbf{Middle} & \textbf{High} \\
\midrule
%%NOVELROWS%%
\bottomrule
\end{tabular}
\begin{tablenotes}\small
\item[a] Source statistic: 38\% judge recycled plastic suitable for food
  packaging; the share is placed on the two highest categories.
\item[b] Four risk statements, reversed, and two safety statements from one
  sample; benefit items are excluded.
\item[c] No usable entry: the only candidate measures perceived benefit.
\end{tablenotes}
\end{threeparttable}
\end{table}

Only food-packaging acceptance has both a base rate and published profile
effects, and only it is projected. For the other four outcomes the evidence
does not support a projection: general acceptance and willingness to pay rest
on heterogeneous items from two studies each, there is no quality measure,
and the single source on safety gives no direction:
respondents disagreed both with risk statements (``carries certain risks'',
``may be contaminated'') and with safety statements (``are considered safe'',
``do not pose any harm''), all with means between 2.0 and 2.2 on a five-point
agreement scale. The data do not support ranking safety above price as a
barrier. Safety-related concerns are prominent in the
source literature \citep{magnier2019,meng2021} and warrant targeted measurement
in Spain; their relative weight is an open empirical question.

%-----------------------------------------------------------------------------
\subsection{The conditional acceptance surface}
\label{sec:surface}
%-----------------------------------------------------------------------------

The object the system returns is a \emph{conditional acceptance surface}: a
mapping from demographic and attitudinal profiles to acceptance distributions,
conditional on the dependency structure learned from national data and on the
priors derived from external evidence. A narrative scenario fixes a few
qualitative futures and an aggregate forecast collapses heterogeneity into one
expected value. The surface keeps segment-level structure and can be queried:
any profile, specified by demographics, environmental concern, recycling
behaviour, or a combination, returns a projected probability. Fully specified
profiles carry an interval; partial queries are answered from the synthetic
population (Section~\ref{sec:step5}). In the
sense of \citet{bas2015}, it treats the future as a space to be explored
through deliberate choices and not as a single outcome to be predicted.

For implementation, the relevant quantity is the composition of demand, not
its mean. A product category bought mainly by low-acceptance profiles carries a
different risk from one bought mainly by high-acceptance profiles, even when
population-wide acceptance is identical. The surface makes that composition
explicit.

The surface is also updatable in a specific sense (R3). When a Spanish survey
measures a target outcome, its counts $\mathbf{n}_{j,c}$ enter
Eq.~\eqref{eq:fusion} as data, shift the posterior toward Spanish values, and
narrow the intervals in proportion to sample size, without re-estimating the
network. It is silent on structural ruptures such as a contamination
scandal, a regulatory reversal, or a material innovation that changes the
choice set. Those remain the province of qualitative foresight, which the
surface complements and does not replace.

%-----------------------------------------------------------------------------
\subsection{Profile spread under the transfer assumptions}
\label{sec:spread}
%-----------------------------------------------------------------------------

Under the central specification ($p_0 = 0.38$, $a_0 = 0.80$), the projected
probability of judging recycled plastic suitable for food packaging is 77\%
(95\%~interval: 58--90\%) for a woman under 30 with higher education who
recycles always or often and is extremely worried about the environment. It is
2\% (0.5--7\%) for a man aged 30 or over without higher education who recycles
sometimes or never and is not at all worried. The difference is
75~percentage points (54--88). The intervals are wide because they carry the
standard errors of five coefficients estimated on 272 respondents.

The two profiles are rare: 18 and 8 of the 2{,}133 CIS respondents match them.
The distribution over the whole sample is more informative for policy. With
CIS weights, the median projected probability is 35\%; a quarter of residents
(26\%) fall below 25\%, and 24\% exceed 50\%. Table~\ref{tab:segments} gives
the projection by segment. The largest contrast is environmental worry (46\%
at the top of the scale against 8\% at the bottom), followed by education
(50\% with higher education against 32\% without), recycling practice (41\%
against 25\%) and gender (45\% for women against 30\% for men).

\begin{table}[htbp]
\centering
\caption{Projected acceptance of recycled plastic in food packaging by
segment of the CIS~3391 sample, weighted with PESOFIN (central specification;
population mean fixed at $p_0 = 0.38$ by calibration).}
\label{tab:segments}
\small
\begin{tabular}{llrrr}
\toprule
\textbf{Dimension} & \textbf{Segment} & $n$ & \textbf{Share} &
  \textbf{Projected acceptance} \\
\midrule
%%SEGROWS%%
\bottomrule
\end{tabular}
\end{table}

Three qualifications apply. The magnitude depends on the transfer assumptions:
the spread between the two profiles ranges from 55 to 83~points for $a_0$
between $0.50$ and $1.00$, and from 62 to 83~points for $p_0$ between $0.25$
and $0.60$ (Appendix~\ref{app:sensitivity}); the ordering of profiles and
segments is unchanged. The segment differences reproduce by construction the
coefficients of one Finnish study, redistributed over the Spanish population
structure. And none of these figures is an observed or validated property of
the Spanish population.

Read as scenarios, the spread separates two situations for a given product
category.

\textbf{Scenario~A (receptive majority).} Environmentally engaged, younger
consumers make up most of the purchasing segment. The recycled-content
requirement meets limited resistance and general communication may suffice.

\textbf{Scenario~B (reluctant majority).} Consumers with low environmental
concern who do not recycle dominate the segment. The same requirement meets
substantial resistance, with consequences for sales and for the perceived
legitimacy of the measure. Nothing in the regulation differs between the two
scenarios; segment composition does.

%-----------------------------------------------------------------------------
\subsection{Sensitivity}
\label{sec:sensitivity}
%-----------------------------------------------------------------------------

Three assumptions are varied (Appendix~\ref{app:sensitivity}). The transfer
factor $a_0$ and the base rate $p_0$ move the spread substantially, as reported
above; the base rate also shifts every level, since the population mean equals
$p_0$ by calibration. The coding of the worry scale is the third: the spread
is 67~points if the coefficient is read as the full contrast between the
extremes and 81~points if it applies per scale point.

Adding the empirical transport error of Section~\ref{sec:transport} (a
random transfer factor, per-coefficient transport noise, and a level shift)
widens every interval. The population mean is no longer fixed at $p_0$:
%%TPOP%%. The most receptive profile is projected at %%TMOST%%, the least
receptive at %%TLEAST%%, and the spread between them at %%TSPREAD%%~points;
segment intervals are given in Table~\ref{tab:plan}. The ordering of profiles survives, and the spread
remains large under the lower bound; the absolute levels are uncertain by
about $\pm 20$ points. These are the intervals that should accompany any use
of the projections.

In the synthetic population, where the outcome is drawn for each of the
%%NWITH%% individuals with complete predictors, %%SHAREACC%% accept, and the
contrasts by gender (%%FEMW%% against %%FEMWO%%), higher education
(%%EDUW%% against %%EDUWO%%) and recycling (%%RECW%% against %%RECWO%%)
reproduce those of the CIS projection in Table~\ref{tab:segments}.

